\appendix

\chapter{Đặc tả Văn phạm Chi tiết của Tập con An toàn (Safe Subset)}
\label{app:safe_subset_grammar}

Phụ lục này cung cấp toàn văn đặc tả ngữ pháp hình thức theo dạng Chuẩn Backus-Naur Mở rộng (Extended Backus-Naur Form -- EBNF) của tập con ngôn ngữ Python QF-LIA được hỗ trợ bởi hệ thống, cùng bảng ánh xạ ngữ nghĩa hình thức tương ứng sang lý thuyết SMT.

\section{Ngữ pháp EBNF}
\begin{lstlisting}[language=Python, caption={Đặc tả EBNF của Safe Subset QF-LIA}]
program      ::= function_def | stmt_list
function_def ::= "def" NAME "(" [param_list] ")" ":" NEWLINE INDENT stmt_list DEDENT
param_list   ::= NAME ("," NAME)*

stmt_list    ::= stmt+
stmt         ::= simple_stmt | compound_stmt

simple_stmt  ::= (assign_stmt | augassign_stmt | assert_stmt | return_stmt | pass_stmt) NEWLINE
compound_stmt::= if_stmt | while_stmt | for_stmt

assign_stmt  ::= NAME "=" expr
augassign_stmt::= NAME ("+=" | "-=") expr
assert_stmt  ::= "assert" bexpr ["," STRING]
return_stmt  ::= "return" [expr]
pass_stmt    ::= "pass"

if_stmt      ::= "if" bexpr ":" suite ("elif" bexpr ":" suite)* ["else" ":" suite]
while_stmt   ::= "while" bexpr ":" suite
for_stmt     ::= "for" NAME "in" "range" "(" range_args ")" ":" suite
suite        ::= NEWLINE INDENT stmt_list DEDENT

range_args   ::= expr | expr "," expr ["," INT_LITERAL]

expr         ::= term (("+" | "-") term)*
term         ::= factor (("*" | "//" | "%") INT_LITERAL)*
factor       ::= ["+" | "-"] atom
atom         ::= INT_LITERAL | NAME | "(" expr ")"

bexpr        ::= bterm ("or" bterm)*
bterm        ::= bfactor ("and" bfactor)*
bfactor      ::= ["not"] bfactor | comp_expr | "(" bexpr ")"
comp_expr    ::= expr comp_op expr
comp_op      ::= "==" | "!=" | "<" | "<=" | ">" | ">="

INT_LITERAL  ::= ["0" | [1-9][0-9]*]
NAME         ::= [a-zA-Z_][a-zA-Z0-9_]*
\end{lstlisting}

\section{Bảng Ánh xạ Ngữ nghĩa sang Lý thuyết SMT (QF-LIA)}
Bảng \ref{tab:semantic_mapping_smt} mô tả quy tắc dịch tương đương từ các cấu trúc cú pháp của Safe Subset sang các toán tử của logic SMT-LIB2 trong Z3.

\begin{table}[H]
\centering
\small
\renewcommand{\arraystretch}{1.2}
\begin{tabular}{|l|l|l|}
\hline
\textbf{Cấu trúc Python} & \textbf{Biểu thức SMT-LIB2 Tương đương} & \textbf{Ngữ nghĩa Toán học} \\
\hline
\texttt{x = e} & \texttt{(= x\_new [[e]])} & Phép gán SSA (Phiên bản mới) \\
\texttt{x += e} & \texttt{(= x\_new (+ x\_old [[e]]))} & Gán kết hợp tăng \\
\texttt{x -= e} & \texttt{(= x\_new (- x\_old [[e]]))} & Gán kết hợp giảm \\
\texttt{e1 + e2} & \texttt{(+ [[e1]] [[e2]])} & Phép cộng số nguyên đại số \\
\texttt{e1 - e2} & \texttt{(- [[e1]] [[e2]])} & Phép trừ số nguyên đại số \\
\texttt{c * x} & \texttt{(* c x)} & Nhân vô hướng với hằng số $c \in \mathbb{Z}$ \\
\texttt{x // c} & \texttt{(div x c)} & Phép chia nguyên (với $c \ne 0$) \\
\texttt{x \% c} & \texttt{(mod x c)} & Phép lấy phần dư (với $c > 0$) \\
\texttt{e1 == e2} & \texttt{(= [[e1]] [[e2]])} & Quan hệ đồng nhất \\
\texttt{e1 != e2} & \texttt{(not (= [[e1]] [[e2]]))} & Quan hệ bất đẳng \\
\texttt{e1 < e2} & \texttt{(< [[e1]] [[e2]])} & Quan hệ thứ tự nghiêm ngặt \\
\texttt{e1 <= e2} & \texttt{(<= [[e1]] [[e2]])} & Quan hệ thứ tự không nghiêm ngặt \\
\texttt{b1 and b2} & \texttt{(and [[b1]] [[b2]])} & Hội logic Boolean \\
\texttt{b1 or b2} & \texttt{(or [[b1]] [[b2]])} & Tuyển logic Boolean \\
\texttt{not b} & \texttt{(not [[b]])} & Phủ định logic Boolean \\
\texttt{if B ...} & \texttt{(ite [[B]] path\_then path\_else)} & Hội tụ luồng điều khiển SSA \\
\hline
\end{tabular}
\caption{Bảng ánh xạ ngữ nghĩa giữa Python Safe Subset và SMT-LIB2}
\label{tab:semantic_mapping_smt}
\end{table}

\chapter{Cấu trúc Dữ liệu}
\label{app:data_schemas}

\section{Khẳng định dùng chung}
Mỗi bản ghi trong \texttt{claim\_pool.json} chứa định danh khẳng định và chương trình, biểu thức Python, số dòng, thông tin neo và nhãn chuẩn. Ví dụ dưới đây trích từ tập dữ liệu.
\begin{lstlisting}[caption={Một bản ghi trong tập khẳng định dùng chung}]
{
  "claim_id": "bench_v1::arithmetic/p00009.py::assert_0",
  "program_id": "arithmetic/p00009.py",
  "kind": "postcondition",
  "text": "_contract_result == x + y",
  "source_line": 4,
  "ast_anchor": "add@Assert:L4:C4#0",
  "precondition": null,
  "expected_verdict": "VERIFIED",
  "gt_confidence": "high",
  "provenance": "assert@L4",
  "validated": {
    "syntax_ok": true,
    "anchor_ok": true
  }
}
\end{lstlisting}

\section{Đặc tả đầu vào của baseline}
Kiểu \texttt{ProgramSpec} có hai trường bắt buộc là \texttt{program\_id} và \texttt{source}. Các trường tùy chọn gồm \texttt{source\_path}, danh sách \texttt{claims}, \texttt{precondition} và từ điển \texttt{metadata}. Danh sách \texttt{claims} sử dụng kiểu dữ liệu định nghĩa trong \texttt{shared/claims.py}.

\chapter{Đặc tả Chi tiết 14 Phương pháp Baseline (A--N)}
\label{app:baselines_specification}

Phụ lục này cung cấp tài liệu kỹ thuật chi tiết về cơ chế vận hành, thuật toán, cờ năng lực và các tham số của toàn bộ 14 baseline đối chứng.

\section{Baseline A: SatLM Direct SMT-LIB2}
\begin{itemize}
    \item \textbf{Cơ chế:} Prompt trực tiếp yêu cầu mô hình đóng vai trò biên dịch viên ngôn ngữ, chuyển toàn bộ logic của hàm Python và thuộc tính kiểm tra thành mã kịch bản SMT-LIB2 hoàn chỉnh chứa các lệnh \texttt{(declare-const ...)}, \texttt{(assert ...)}, và lệnh kiểm tra phủ định thuộc tính \texttt{(check-sat)}.
    \item \textbf{Bộ giải:} Script SMT-LIB2 được chuyển thẳng cho Z3 Solver. Nếu kết quả là \texttt{unsat} $\implies$ \texttt{VERIFIED}; nếu là \texttt{sat} $\implies$ \texttt{COUNTEREXAMPLE}.
    \item \textbf{Đặc tính:} Không có AST trung gian, không có truy vấn Reachability, không sinh input cụ thể để replay CPython.
\end{itemize}

\section{Baseline B: Execution Grounding}
\begin{itemize}
    \item \textbf{Cơ chế:} Hệ thống tự động khởi chạy chương trình trên một lưới gồm 10 vector đầu vào cụ thể (bao gồm các giá trị biên $-1, 0, 1$). Ghi nhận vết thực thi động (dòng lệnh đã chạy và biến sau mỗi dòng).
    \item \textbf{Prompting:} Vết thực thi được định dạng thành bảng văn bản và đưa vào prompt: \textit{``Dựa trên các vết thực thi đã quan sát được ở trên, hãy xác định xem câu lệnh assert có luôn đúng hay không.''}
    \item \textbf{Hỗ trợ Replay:} Có. Mọi phản ví dụ được kiểm tra chéo ngay trên runtime.

\end{itemize}

\section{Baseline C: Structural Reasoning (CFG + Python k-induction)}
\begin{itemize}
    \item \textbf{Cơ chế:} Trích xuất bộ khung điều khiển (CFG) từ AST. Triển khai thuật toán $k$-induction hoàn toàn trong môi trường Python bằng cách duỗi các bước lặp tượng trưng và đánh giá biểu thức mà không gọi Z3.
    \item \textbf{Ưu nhược điểm:} Độc lập hoàn toàn với solver hình thức bên ngoài, nhưng bị giới hạn về khả năng giải quyết các hệ bất phương trình đa biến phức tạp.
\end{itemize}

\section{Baseline D: LLM-as-a-Judge}
\begin{itemize}
    \item \textbf{Cơ chế:} Đại diện cho phương pháp đánh giá AI zero-shot phổ biến. LLM nhận mã nguồn và claim, sinh chuỗi lý luận CoT bằng ngôn ngữ tự nhiên tự do, tự phân tích các trường hợp biên, và đưa ra câu trả lời theo cấu trúc: \texttt{VERDICT: [VERIFIED | COUNTEREXAMPLE]}.
    \item \textbf{Phân tích Phản ví dụ:} Nếu mô hình tuyên bố \texttt{COUNTEREXAMPLE}, hệ thống sử dụng biểu thức chính quy (Regex) để trích xuất các giá trị biến do mô hình đề xuất và kích hoạt CPython replay.
\end{itemize}

\section{Baseline E: Neuro-Symbolic AutoSpec}
\begin{itemize}
    \item \textbf{Cơ chế:} Mô phỏng công trình AutoSpec \cite{wen2024autospec}. LLM đề xuất một bộ đặc tả đầy đủ (precondition, postcondition, invariant). Bộ kiểm tra hình thức tĩnh đối chiếu đặc tả với mã nguồn. Nếu phát hiện vi phạm hoặc đặc tả mâu thuẫn, thông báo lỗi được phản hồi lại để LLM tinh chỉnh trong tối đa 3 vòng lặp.
\end{itemize}

\section{Baseline F: Agentic Multi-Step (ReAct-style)}
\begin{itemize}
    \item \textbf{Cơ chế:} Tác tử hoạt động theo chu trình tự hành ReAct:
    \begin{enumerate}
        \item \texttt{PLAN:} Lập kế hoạch kiểm chứng các trường hợp rẽ nhánh.
        \item \texttt{ANALYZE:} Phân tích logic chuyển trạng thái của biến.
        \item \texttt{TEST:} Tác tử chủ động gọi công cụ Tool để chạy thử chương trình trên CPython với một vector đầu vào nghi ngờ.
        \item \texttt{VERIFY:} Dựa trên kết quả chạy thật, xác nhận hoặc bác bỏ giả thuyết.
        \item \texttt{REFINE:} Tinh chỉnh kết luận cuối cùng.
    \end{enumerate}
\end{itemize}

\section{Baseline G: Self-Consistency}
\begin{itemize}
    \item \textbf{Cơ chế:} Khai thác kỹ thuật Self-Consistency \cite{wang2022self}. Khởi tạo 5 phiên lấy mẫu CoT độc lập từ Gemini với nhiệt độ $T = 0.7$. Trích xuất phán quyết từ cả 5 phiên và áp dụng luật bỏ phiếu đa số tuyệt đối (Majority Voting). Nếu tỷ lệ phiếu hòa (ví dụ 2 VERIFIED, 2 COUNTEREXAMPLE, 1 UNKNOWN), hệ thống trả về \texttt{UNKNOWN\_TIMEOUT}.
\end{itemize}

\section{Baseline H: Self-Refine}
\begin{itemize}
    \item \textbf{Cơ chế:} Khai thác kỹ thuật Self-Refine \cite{madaan2024self}. Mô hình thực hiện quy trình tự phản biện nội tâm qua 3 vòng lặp liên tiếp:
    \begin{itemize}
        \item Vòng 1: Đưa ra phán quyết ban đầu và lời giải thích.
        \item Vòng 2: Tự đóng vai một chuyên gia phản biện khắt khe, tìm kiếm các lỗi ngụy biện hoặc ca biên bị bỏ sót trong lời giải của Vòng 1.
        \item Vòng 3: Dựa trên lời tự phê bình, cập nhật lại phán quyết và lời giải thích cuối cùng.
    \end{itemize}
\end{itemize}

\section{Baseline I: LLM-Guided Property-Based Testing (PBT)}
\begin{itemize}
    \item \textbf{Cơ chế:} LLM phân tích mã nguồn và tự động sinh mã kiểm thử dựa trên thuộc tính sử dụng thư viện Hypothesis của Python. Mô hình định nghĩa các chiến lược sinh dữ liệu (Strategies như \texttt{st.integers(min\_value=0, max\_value=100)}). Thư viện Hypothesis chạy 200 ca kiểm thử ngẫu nhiên có thu nhỏ (shrinking) để tìm phản ví dụ tối thiểu.
\end{itemize}

\section{Baseline J: Deterministic Fuzzing (0 Token, Offline)}
\begin{itemize}
    \item \textbf{Cơ chế:} Phương pháp kiểm thử tự động cổ điển, tiêu thụ đúng 0 token. Bộ fuzzing sinh 256 vector thử nghiệm dựa trên phân tích giá trị biên (Boundary Value Analysis): tập các số cực tiểu/cực đại, $0, 1, -1, 2, -2$, các lũy thừa của 2 ($2^8, 2^{16}, 2^{31}-1$), và các số giả ngẫu nhiên từ seed 42. Chạy toàn bộ trên CPython.
\end{itemize}

\section{Baseline K: Mutation Testing (0 Token, Offline)}
\begin{itemize}
    \item \textbf{Cơ chế:} Tạo ra 50 biến thể đột biến cú pháp (mutants) của chương trình gốc:
    \begin{itemize}
        \item Hoán đổi toán tử quan hệ: $<$ thành $\le$, $==$ thành $!=$, v.v.
        \item Hoán đổi toán tử số học: $+$ thành $-$, $*$ thành $/$.
        \item Dịch chuyển hằng số: $c$ thành $c + 1$ hoặc $c - 1$.
    \end{itemize}
    Chạy claim đối chiếu với các mutant để đo điểm đột biến (Mutation Score).
\end{itemize}

\section{Baseline L: Pure Symbolic Execution (0 Token, Offline)}
\begin{itemize}
    \item \textbf{Cơ chế:} Thực thi ký hiệu thuần túy không dùng LLM dựa trên Z3 và thuật toán PDR/Spacer. Liệt kê toàn bộ các đường đi khả đạt tới câu lệnh \texttt{assert}. Hỗ trợ cờ \texttt{supports\_reachability = True}.
\end{itemize}

\section{Baseline M: Hoare Triple Natural Language Verification}
\begin{itemize}
    \item \textbf{Cơ chế:} LLM được hướng dẫn phân rã toàn bộ hàm thành chuỗi các bộ ba Hoare $\{P\} \; C \; \{Q\}$ bằng văn bản tự nhiên, phân tích tính đúng đắn của từng bộ ba theo các tiên đề Hoare và tổng hợp kết luận.
\end{itemize}

\section{Baseline N: Invariant Synthesis with SMT}
\begin{itemize}
    \item \textbf{Cơ chế:} LLM chỉ làm nhiệm vụ duy nhất là đề xuất từ 1 đến 3 ứng viên bất biến vòng lặp. Z3 solver đóng vai trò thẩm định viên độc lập, kiểm tra Base Case và Inductive Step. Nếu bất biến được chứng minh, Z3 dùng nó để xác minh assertion.
\end{itemize}

\chapter{Mẫu Nhắc lệnh (Prompt Templates)}
\label{app:prompt_templates}

\section{Sinh khẳng định từ mã nguồn}
Mô hình nhận mã nguồn đánh số dòng và được yêu cầu trả về một đối tượng JSON gồm \texttt{cot\_trace} và \texttt{claims}. Mỗi đề xuất chỉ rõ một số dòng và một biểu thức Boolean Python. Biểu thức được kiểm tra ngay trước khi câu lệnh ở dòng đó thực thi.

\section{Phản hồi hồi cứu}
Prompt cung cấp mã nguồn, diễn giải trước đó, claim ban đầu và phản ví dụ đã được CPython xác nhận. Ngữ cảnh gồm đầu vào, dòng vi phạm và trạng thái ghi nhận tại đó. Phản hồi có thể gồm \texttt{refined\_cot}, \texttt{refined\_claim} với hai trường \texttt{line} và \texttt{expression}, \texttt{explanation} và \texttt{accept\_ce}.

\chapter{Các Ví dụ Minh họa}
\label{app:case_studies}

Phụ lục này dùng ba chương trình nhỏ để minh họa cách lập nghĩa vụ chứng minh, kiểm tra bất biến và tái hiện phản ví dụ. Các bước dưới đây được diễn giải từ mã nguồn và công thức của từng ví dụ.

\section{Ca Nghiên cứu 1: Kiểm chứng Vòng lặp với Bất biến Quy nạp}
\begin{lstlisting}[language=Python, caption={Mã nguồn bài toán loop\_sum\_safe}]
def loop_sum(n: int) -> int:
    # Precondition: n >= 0
    s = 0
    i = 0
    while i < n:
        s = s + 1
        i = i + 1
    assert s == n
    return s
\end{lstlisting}

\begin{itemize}
    \item \textbf{Cổng neo AST:} Khẳng định \texttt{s == n} gắn với câu lệnh \texttt{assert} ở dòng 8 và được xét trên trạng thái ngay trước khi câu lệnh chạy.
    \item \textbf{Tổng hợp bất biến:} Xét ứng viên bất biến sau.
    \begin{equation*}
    I(s, i, n) \equiv (s == i) \land (0 \le i \le n)
    \end{equation*}
    \item \textbf{Kiểm chứng Z3:}
    \begin{enumerate}
        \item \textit{Khởi tạo:} $(n \ge 0 \land s = 0 \land i = 0) \implies (s == i \land 0 \le i \le n)$ $\implies$ \texttt{UNSAT} cho phủ định $\implies$ Đúng.
        \item \textit{Bảo toàn:} $(s == i \land 0 \le i \le n \land i < n \land s' = s + 1 \land i' = i + 1) \implies (s' == i' \land 0 \le i' \le n)$ $\implies$ \texttt{UNSAT} cho phủ định $\implies$ Đúng.
        \item \textit{Hậu điều kiện:} Khi vòng lặp kết thúc, $\neg (i < n) \implies i \ge n$. Kết hợp với $i \le n$ từ bất biến suy ra $i == n$. Vì $s == i$, suy ra $s == n$.
    \end{enumerate}
    \item \textbf{Phán quyết Cuối cùng:} \texttt{VERIFIED}.
\end{itemize}

\section{Ca Nghiên cứu 2: Phát hiện Phản ví dụ và Tự Sửa sai Thành công}
\begin{lstlisting}[language=Python, caption={Mã nguồn bài toán bước nhảy step\_two\_buggy}]
def step_two(n: int) -> int:
    # Precondition: n >= 0
    x = 0
    while x < n:
        x = x + 2
    assert x == n
    return x
\end{lstlisting}

\begin{itemize}
    \item \textbf{Khẳng định ban đầu:} \texttt{x == n} tại dòng 6.
    \item \textbf{Z3 Validity Query:} Solver tìm được nghiệm phản ví dụ: $n = 1$.
    \begin{itemize}
        \item Khởi tạo: $x = 0$. Vòng lặp kiểm tra: $0 < 1$ (True).
        \item Thân vòng lặp: $x = 0 + 2 = 2$.
        \item Vòng lặp kiểm tra tiếp: $2 < 1$ (False) $\implies$ Thoát vòng lặp.
        \item Tại dòng assert: kiểm tra $x == n \iff 2 == 1 \implies$ Sai!
    \end{itemize}
    \item \textbf{CPython replay:} Với $n = 1$, câu lệnh \texttt{assert} tại dòng 6 nhận $x = 2$ và phát sinh \texttt{AssertionError}.
    \item \textbf{Phản hồi hồi cứu:} Trạng thái $n = 1, x = 2$ cung cấp ngữ cảnh để đề xuất một khẳng định phụ thuộc vào tính chẵn lẻ của $n$.
    \item \textbf{Khẳng định được tinh chỉnh:} Xét tính chất sau tại vị trí neo.
    \begin{equation*}
    \text{Revised Claim:} \quad (n \% 2 == 0 \implies x == n) \land (n \% 2 == 1 \implies x == n + 1)
    \end{equation*}
    \item \textbf{Kiểm chứng lại:} Nếu chứng minh được khẳng định tinh chỉnh, phán quyết áp dụng cho chính khẳng định đó. Với $n = 1$, câu lệnh \texttt{assert x == n} gốc vẫn vi phạm.
\end{itemize}

\section{Ca Nghiên cứu 3: Ngăn chặn Chứng minh Rỗng trên Nhánh Chết}
\begin{lstlisting}[language=Python, caption={Mã nguồn chứa nhánh chết unfeasible\_branch}]
def dead_code_demo(a: int):
    # Precondition: True
    if a > 10:
        if a < 5:
            # Nhanh chet: a > 10 and a < 5 la vo nghiem
            assert 1 == 0
\end{lstlisting}

\begin{itemize}
    \item \textbf{Kiểm tra tính hợp lệ:} Truy vấn $(a > 10 \land a < 5) \land \neg (1 == 0)$ vô nghiệm do điều kiện đường đi mâu thuẫn.
    \item \textbf{Proposed Pipeline:} Giao thức Hai Truy vấn thực thi Truy vấn 1 (Reachability):
    \begin{equation*}
    \Phi_{\text{reach}} \equiv \exists a. \ (a > 10 \land a < 5) \implies \texttt{UNSAT}
    \end{equation*}
    Truy vấn khả đạt cho phép gắn nhãn \texttt{UNREACHABLE} cho khẳng định tại vị trí này.
\end{itemize}
